%% CRS minimal art-paper skeleton — SIGGRAPH Asia Art Papers (acmart).
%% Demonstrates the canonical CRS output path: acmart sigconf + ACM Reference
%% Format, Practice-Based Art Paper structure (Pattern 1), and the genre-specific
%% artwork citation (@misc with venue+date as the L3 locator).
%%
%% Build:  latexmk -pdf paper.tex      (or: pdflatex; bibtex; pdflatex x2)
%% Requires the ACM `acmart` package (CTAN / texlive-publishers). Verify the
%% exact class option for the current SIGGRAPH Asia Art Papers track CFP;
%% `sigconf` is the CRS default.
\documentclass[sigconf,nonacm]{acmart}

\settopmatter{printacmref=false}
\setcopyright{none}
\acmConference[SA Art Papers '26]{SIGGRAPH Asia Art Papers}{2026}{(verify venue/date against the current CFP)}

\begin{document}

\title{A Practice-Based Art Paper Skeleton}
\subtitle{Canonical CRS acmart output (template, not a submission)}

\author{Artist Name}
\affiliation{\institution{Studio / Institution}\country{Country}}
\email{artist@example.org}

\begin{abstract}
This skeleton exercises the Creative Research Skills (CRS) output path: an
\texttt{acmart} \texttt{sigconf} document, ACM Reference Format citations, the
Practice-Based Art Paper structure, and a genre-specific artwork citation. In a
real paper this abstract states the work, its provocation, the essential making
gesture, the situated insight, and the contribution to the art-and-technology
discourse, in roughly 120--200 words.
\end{abstract}

\keywords{practice-based research, media art, generative art}

\maketitle

\section{Introduction / Context}
The work in one paragraph; its artistic and conceptual context; the provocation
it pursues; and the contribution this paper makes. Precedent work is positioned
here, e.g. a precedent artwork~\cite{precedentArtwork} and relevant
theory~\cite{leonardoArticle}.

\section{Conceptual Framework}
The theoretical grounding, the relationship to precedent artworks and artists,
and the key concepts and their definitions. The \emph{artwork is the primary
evidence}; the framework is what the work argues through.

\section{The Work}
Form, materials, media, scale, duration; what the audience perceives or does;
authorship and collaboration credit.

\section{Realization / Methods of Making}
The technical approach, the process and its iterations, and the tools and
dependencies. Technical claims (``real-time'', ``autonomous'') are anchored to
this section, never asserted without realization evidence~\cite{siggraphWork}.

\section{Reflection / Discussion}
What the making revealed (situated insight); exhibition and reception with
observable anchors (named venue + date, not ``audiences were moved''); relation
back to the framework; limitations and open questions.

\section{Conclusion}
What the work opens up, and the future trajectory.

\begin{acks}
\textbf{Use of AI Tools.} In producing \emph{this paper}, the authors used
AI-assisted tooling (Creative Research Skills) for drafting and citation
formatting. [If applicable:] as part of \emph{the artwork itself}, generative AI
was used as a medium, described in the realization section above. The authors
reviewed all content and take full responsibility for it. No AI system is an
author of this work. (Verify wording against the current ACM / SIGGRAPH Asia
policy.)
\end{acks}

\bibliographystyle{ACM-Reference-Format}
\bibliography{refs}

\end{document}
